Il presente report descrive l'intero processo di progettazione, sviluppo e valutazione di \textbf{Vivo}, un'applicazione mobile sviluppata in Dart utilizzando il framework Flutter.
Il progetto affronta una problematica attuale e molto sentita, ossia la salute e il benessere mentale, con lo scopo specifico di aiutare le persone a interrompere i meccanismi dell'ansia e degli attacchi di panico.
L'obiettivo dell'applicazione è fornire strumenti e risorse per la gestione dell'ansia, offrendo un supporto pratico e immediato agli utenti.
Questo documento illustra la realizzazione dell'applicazione seguendo, passo dopo passo, i principi dello \textit{User-Centered Design} (UCD), al fine di garantire un'esperienza utente ottimale e soddisfacente.

\section{Organizzazione del report}
\label{sec:organizzazione-del-report}

Questo report è strutturato in diverse sezioni, ciascuna delle quali affronta un aspetto specifico del progetto.
Il percorso metodologico è suddiviso in tre step sequenziali:

\begin{itemize}
    \item \textbf{Step 1:} Envisioning. Questa fase è dedicata a definire l'identità dell'applicazione, a posizionarla sul mercato rispetto ai competitor e a delineare e stabilire le priorità dei requisiti principali che rispondono alle esigenze degli utenti target.
    \item \textbf{Step 2:} Prototipazione e valutazione. Questa fase si concentra sull'esplorazione e sulla definizione dell'interfaccia utente (UI), partendo da prototipi a bassa fedeltà disegnati a mano (\textit{sketching}), fino ad arrivare a prototipi a media fedeltà realizzati con strumenti digitali come Figma (\textit{wireframing}). Questi ultimi vengono infine sottoposti a valutazioni e test da parte di utenti reali.
    \item \textbf{Step 3:} Sviluppo. Questa fase consiste nella traduzione del design in un'applicazione mobile vera e propria, accompagnata da una valutazione critica delle tecnologie utilizzate e delle variazioni apportate rispetto al prototipo iniziale.
\end{itemize}

Ai tre step corrispondono i capitoli che seguono.
Il capitolo 2 copre il primo step, con il concept statement, l'analisi dei competitor, le personas e i requisiti.
I capitoli 3, 4 e 5 coprono il secondo, seguendo l'interfaccia dagli sketch a mano libera ai wireframe digitali fino alla valutazione con gli utenti.
Il capitolo 6 copre il terzo, con le tecnologie adottate, le schermate dell'applicazione realizzata e le differenze rispetto al prototipo.
Il capitolo 7 conclude il documento con un bilancio del lavoro svolto e propone i possibili sviluppi futuri dell'applicazione.

In sintesi, il presente documento intende dimostrare l'efficacia di un design empatico e mirato.
L'obiettivo è evidenziare come un'interfaccia intuitiva possa trasformarsi in uno strumento di supporto concreto, capace di assistere tempestivamente l'utente nella gestione quotidiana dell'ansia e nei momenti di maggiore vulnerabilità.
